\section{The small algebra moves we will use}
An expression such as $2k+1$ is built from operations you already know. To evaluate it at $k=3$, replace every occurrence of $k$ by $3$: $2k+1$ becomes $2\cdot3+1=7$. This replacement is called \emph{substitution}. Substitution does not say that $k$ always equals three; it checks one chosen input.

Sometimes we give the rule itself a name. Writing $f(n)=n+1$ means that the rule named $f$ adds one to its input $n$; for example, $f(3)=4$. Read $f(n)$ as ``the output of $f$ at input $n$,'' not as $f$ multiplied by $n$. Chapter~3 develops functions fully. For now this notation lets us refer to a formula and its values without rewriting the formula each time.

To rewrite an expression without changing its value, use an arithmetic rule in either direction. Distributing gives $2(k+1)=2k+2$. Reading that equality in reverse gives $2k+2=2(k+1)$; this is called \emph{factoring}. In a proof, factoring is often useful because it reveals a common multiplier.

For a slightly longer example, apply distribution one bracket at a time:
\begin{align*}
 (2k+1)(2\ell+1)
 &=2k(2\ell+1)+1(2\ell+1)\\
 &=4k\ell+2k+2\ell+1.
\end{align*}
The first line gives each term of the first bracket a copy of the second bracket. The second line multiplies out those two copies. The letter $\ell$, read ``ell,'' is a different name from $k$, not the digit one. We need two names because the inputs may differ.

There is also a difference between an identity and an equation to solve. The equality $2(k+1)=2k+2$ holds for every integer $k$; it is an \emph{identity}. The equation $2k+1=7$ asks which values of $k$ make it true. Subtracting one from both sides gives $2k=6$; dividing by two gives $k=3$. Substitution into the original equation checks that $2\cdot3+1=7$.

A power records repeated multiplication: $2^3=2\cdot2\cdot2=8$. The exponent three counts how many factors are multiplied; it does not count additions, so $2^3$ is not $2+2+2$.

Two rules follow directly from counting factors. If we write a row of two factors of $2$ next to a row of three, we get a row of five:
\[
2^2\cdot2^3=(2\cdot2)(2\cdot2\cdot2)=2^5.
\]
In general $a^na^m=a^{n+m}$ for positive integers $n,m$. Division cancels factors instead: $2^5/2^2=(2\cdot2\cdot2\cdot2\cdot2)/(2\cdot2)=2^3$. In general $a^n/a^m=a^{n-m}$ when $n>m$ and $a\ne0$.

We extend the notation to the exponents zero and negative integers so that these two rules keep working. Dividing $a^n$ by itself gives $1$, and the subtraction rule would give $a^{n-n}=a^0$; so we agree that $a^0=1$ whenever $a\ne0$. Similarly, $a^0/a^n$ should be $a^{0-n}=a^{-n}$, so we agree that $a^{-n}=1/a^n$ for a positive integer $n$ and nonzero $a$. For example, $2^{-3}=1/2^3=1/8$. With these agreements, $a^na^m=a^{n+m}$ and $a^n/a^m=a^{n-m}$ hold for all integers $n,m$ whenever $a\ne0$. We will use these rules with actual numbers, especially powers such as $2^{-n}$, from Chapter~5 onward.

\pauseq{Why is $(a+b)^2$ not generally $a^2+b^2$?}{The square means $(a+b)(a+b)$. Distribution gives $a^2+ab+ba+b^2=a^2+2ab+b^2$. For $a=1,b=2$, the two proposed values are nine and five. The missing cross terms account for the difference.}
